\section{Ask what follows from the complete family}
\label{sec:family-outcome-queries}

The retained histories answer a financial question that a single simulated
path cannot settle: which amounts could occur, and which amounts must occur
under the supplied assumptions? In the preceding example, inventory of
65 dollars after settlement is possible. Inventory of 60 dollars after the
whole-lot sale is necessary. The distinction depends on all admissible paths,
including paths that a first witness does not select.

\subsection{Separate a question from an observation}

An observation of inventory of 65 dollars restricts the historical family. Asking whether
inventory of 65 dollars is possible leaves that family unchanged. The query therefore
carries no additional source assertion. It retains the original opening,
event order, branch conditions, funding restrictions, observations and premises.
It also retains unresolved obligations concerning the actual company.

Write \(\mathcal H\) for this complete conditional family and \(h\) for one
history within it. At a chosen prefix \(t\in\{1,2\}\), let
\(A_{a,t}(h)\) be the signed balance of account \(a\) after that event.
The query asks about the scalar amount
\begin{equation}
\label{eq:family-outcome-amount}
 f_t(h)=\sum_{a\in\mathcal A}c_a A_{a,t}(h),
 \qquad h\in\mathcal H.
\end{equation}
Here \(\mathcal A\) is the set of distinct queried accounts, and each
\(c_a\) is an exact rational, dimensionless coefficient. Every account must
have the same full monetary type at that prefix: entity, legal book, date,
currency, dimensions and scale. A threshold has that same monetary type.
This is a same-date account sum, not permission to add differently dated or
differently denominated amounts. The retained profit query applies the signed
account convention to revenue, cost of goods sold and selling expense.

\path{SequentialOutcomeQuery/1} contains the full \path{family}, the chosen
\path{prefix}, the complete \path{type}, the account \path{terms}, an
\path{operator} and a rational \path{threshold}, alongside its \path{profile}
identifier. The five comparisons are equality, strict less or greater, and
weak less or greater. Arbitrary callbacks and compound expressions are outside
this first interface.

\subsection{Distinguish possible, necessary and empty}

Let \(Q(h)\) denote the selected comparison of \(f_t(h)\) with its threshold.
The implementation asks two existence questions on the unchanged family:
\begin{equation}
\label{eq:family-outcome-partition}
 \mathcal H_Q=\{h\in\mathcal H:Q(h)\},\qquad
 \mathcal H_{\neg Q}=\{h\in\mathcal H:\neg Q(h)\}.
\end{equation}
A checked member of \(\mathcal H_Q\) establishes possibility. A checked
member of \(\mathcal H_{\neg Q}\) is a counterexample to necessity. If both
exist, the comparison is contingent. If the first exists and the second is
proved empty, it is necessary within this family. If the first is proved empty
and the second exists, it is impossible. When both sets are empty, the answer
is \path{empty_family}; necessity is not applicable. In particular, an
unfunded settlement does not make every financial assertion meaningfully true.

For inventory of 65 dollars after settlement, a path with \(V=5\) is a satisfying
witness and a path with \(V=20\) is a counterexample. For terminal inventory of 60 dollars,
equation~\eqref{eq:native-sequence-balances} supplies the invariant, while
the complete branch check proves that no counterexample meets the original
conditions. Inventory of 64 dollars after settlement is impossible in this example.
Adding the observation inventory of 65 dollars first would instead make that amount
necessary in the smaller, observed family. These are different questions with
different conditioning information.

Every witness replays both original native transitions. The verifier computes
the account sum independently with exact rational arithmetic and checks which
side of the comparison it satisfies. An empty region requires complete
infeasibility evidence and a separately checked unsatisfiable constraint core.
The checker accepts a different valid witness or a different valid core;
reproducing the solver's first choice is not the criterion.

Search is divided into elementary regions. For equality, the regions are
equal, less and greater; for each inequality, its complement uses the exact
opposite strict or weak comparison. One shared cell-count and elapsed-time budget spans all
regions of all parent branches. An unvisited region cannot prove emptiness.
A found counterexample can already refute necessity, even if the search is
otherwise incomplete; the result preserves that fact while leaving unproved
conclusions unknown. The number of witnesses supplies no probability.

\subsection{Bound the amount without selecting a history}

The closed \path{SequentialOutcomeRange/1} request contains its \path{profile},
\path{family}, \path{prefix}, \path{type} and account \path{terms}. It asks
for the tight lower and upper endpoints instead of a comparison. For a
nonempty family these are
\begin{equation}
\label{eq:family-outcome-envelope}
 L=\inf_{h\in\mathcal H} f_t(h),\qquad
 U=\sup_{h\in\mathcal H} f_t(h).
\end{equation}
An endpoint is reported as attained only if an original admissible history
reaches it. Thus an open interval has valid finite endpoints without supplying
executable states at those endpoints.

The supported fragment has finitely many affine branches with explicitly
finite numeric input domains. Each feasible branch is a convex region defined
by rational linear equalities and strict or weak inequalities. Within a
nonempty such branch, a closure point can be approached from an original
feasible point: their convex combinations satisfy every strict inequality
until the limiting closure point is reached. Optimizing over the closure can
therefore locate its infimum or supremum. Certification additionally checks
that the endpoint is feasible in that closure, no original value is better,
and the original strict constraints either admit or exclude endpoint equality.
An attained endpoint must also pass complete native replay.

The global lower endpoint is the smallest certified branch infimum; the upper
is the largest certified branch supremum. This aggregation requires every
original branch to be either checked infeasible or certified at both endpoints.
A partial search returns local branch evidence with no certified global bounds.
An empty family has no bounds. An unbounded answer is invalid in this explicitly
finite fragment, rather than an authorized infinity result.

The amounts in the example follow directly from the shared carrying value.
With returned quantity fixed at two, \(V=\min(20,N)\). The supplied recoverable
value domain is \(5\leq N\leq30\), so it permits both minimum branches and the carrying
envelope is \([5,20]\). Substituting into
equations~\eqref{eq:native-sequence-balances} and
\eqref{eq:native-sequence-profit} gives
\begin{align}
\label{eq:family-outcome-example}
 I_1&=60+V\in[65,80], & I_2&=60,\\
 \pi_1&=V-24\in[-19,-4], & \pi_2&=P-24.
\end{align}
At proceeds \(P=40\), cumulative terminal profit is 16 dollars before tax
closing. If proceeds instead range over \([0,100]\), its envelope is
\([-24,76]\). Opening inventory of 60 dollars is included in both inventory queries.
These equalities use the original signed postings; selecting two convenient
witnesses would not establish the bounds.

If recoverable value is restricted to \(5<N<20\), then \(V=N\), and the
intermediate inventory envelope becomes \((65,80)\): neither endpoint is
attained. In the original two-branch family, both branches support the upper
bound of 80 dollars, but only the branch with \(V=20\) attains it. The other approaches 80
from below. Accordingly, \path{supporting_branches} and
\path{attaining_branches} are different result fields. Global attainment
requires at least one attaining branch at the global endpoint.

Even a tight outer envelope need not describe every feasible amount. A union
of branches could permit only 65 and 80, leaving the intervening values
unattainable while retaining the same endpoints. The range interface therefore
reports an envelope and endpoint membership, not a general proof of interior
feasibility. A predicate query can separately test a particular interior amount.

The reference definitions and saved complete requests/results are in
\path{data/accounting-master/}:
\begin{quote}\small
\path{sequential-outcome-query-20261024}\\
\path{sequential-outcome-range-20261024}
\end{quote}
Both interfaces preserve missing premises and unresolved target obligations.
Where those obligations remain, an answer concerns the candidate family and
the actual outcome stays unknown. Their states are hypothetical and uncommitted;
source admission, accounting admission and method eligibility remain false.
General histories, nonlinear or integer-lot methods, actual-company state
reconstruction and independent implementations remain outside these checked
references.
